\documentclass[12 pt, letterpaper]{hmcpset}
\usepackage{hw-preamble}

\name{Jayden Thadani}
\class{MATH 131 --- Mathematical Analysis I}
\assignment{Homework 2}
\duedate{14th September 2024}

\begin{document}

\begin{problem}[1.]
	Let $x = a_{0}.a_{1} \cdots$, $y = b_{0}.b_{1} \cdots$ be positive real numbers with $a_{k} \le b_{k}$ for all $k = 0, 1, \dots$. Prove that $z = \sup \{ b_{0}.b_{1} \cdots b_{k} - a_{0}.a_{1} \cdots a_{k} \mid k \in \mathbb{N} \}$ satisfies $x + z = y$.
\end{problem}

\pagebreak

\begin{problem}[2.]
	If $x$ is a positive real number, prove that there exists a positive real number $y$ such that $y^{2} = x$.
\end{problem}

\pagebreak

\begin{problem}[3.]
	Let $p$ and $q$ be relatively prime positive integers and $x \in \mathbb{R}_{+}$ be such that $xq = p$. Prove that $x \in \mathbb{D}$ if and only if the only prime factors of $q$ are $2$ and $5$.
\end{problem}

\pagebreak

\begin{problem}[4 (Rudin chapter 2, problem 2).]
	A complex number $z$ is said to be \textit{algebraic} if there are integers $a_{0}, \dots, a_{n}$ not all zero, such that
	\[
		a_{0} z^{n} + a_{1} z^{n - 1} + \dots + a_{n - 1} z + a_{n} = 0.
	\]
	Prove that the set of all algebraic numbers is countable.	
\end{problem}

\pagebreak

\begin{problem}[5.]
	\begin{enumerate}
		\item Let $x, y$ be positive real numbers with periodic decimal expansions $0.0a_{2} \cdots a_{k} 0 a_{2} \cdots a_{k} \cdots$ and $0.0 b_{2} \cdots b_{k} 0 b_{2} \cdots b_{k} \cdots$ respectively. Prove that the decimal expansion of $x + y$ is also periodic.
		\item Let $x, y$ be positive real numbers with periodic decimal expansions $0.0 a_{2} \cdots a_{k} 0 a_{2} \cdots a_{k} \cdots$ and $0.0 b_{2} \cdots b_{k} 0 b_{2} \cdots b_{k} \cdots$ respectively. Prove that the decimal expansion of $x + y$, $c_{0}.c_{1} \cdots$ satisfies $c_{i} = c_{i + kj}$.
	\end{enumerate}
\end{problem}

\pagebreak

\begin{problem}[6.]
	Let $x < y$ be positive real numbers. Prove that $\{ z \in \mathbb{R} \setminus \mathbb{Q} \mid x < z < y \}$ is uncountable.
\end{problem}

\end{document}
